\chapter{Marco Teórico}
\label{cap:marco_teorico}



Con el fin de optimizar el flujo operativo dentro de la Defensoría de Derechos
Politécnicos (DDP), se implementará un sistema de clasificación para determinar
la competencia o incompetencia de una queja. Este sistema brindará apoyo
directamente en la etapa de primer contacto; con ello se busca transformar el
análisis normativo en un modelo computacional capaz de identificar la competencia
de manera consistente y eficiente. Para lograrlo, el sistema se apoya en cuatro
pilares conceptuales que se desarrollan en este capítulo: el marco legal que
define el problema de clasificación, los fundamentos del Procesamiento de
Lenguaje Natural que permiten interpretar la narrativa de las quejas, las técnicas
de representación vectorial que traducen ese lenguaje en señales numéricas, y los
modelos de aprendizaje automático que aprenden a distinguir entre quejas
competentes e incompetentes.

% ─────────────────────────────────────────────────────────────────────────────
\section{Problema de clasificación en el contexto de la DDP}
% ─────────────────────────────────────────────────────────────────────────────

El problema central de este proyecto consiste en automatizar la decisión de
admisibilidad de una queja ante la Defensoría de Derechos Politécnicos, decisión
que actualmente depende del criterio manual del personal de primer contacto y
genera inconsistencias operativas. Este problema parte del marco legal presentado
en el Acuerdo publicado en la Gaceta Politécnica, Número Extraordinario 622,
donde se establece a la DDP como el órgano encargado de promover, proteger y
defender los derechos de la comunidad frente a actos u omisiones de las
autoridades del Instituto. Por lo tanto, comprender los límites normativos de la
competencia institucional es el primer paso para diseñar un clasificador que
replique la lógica jurídica con precisión.

En el Artículo 3° del mismo Acuerdo se establece que la Defensoría es competente
para conocer quejas relacionadas con presuntas violaciones a los derechos de la
comunidad politécnica; sin embargo, el Artículo 4° delimita los márgenes en las
materias en que la Defensoría \textbf{no} es competente:

\begin{enumerate}
  \item Afectaciones de carácter colectivo.
  \item Resoluciones disciplinarias.
  \item Derechos de naturaleza laboral.
  \item Evaluaciones académicas de profesores, comisiones dictaminadoras o
        consejos técnicos consultivos escolares, salvo que en dicha evaluación
        se violen notoriamente los derechos de los interesados.
  \item Procedimientos y resoluciones instaurados por el Órgano Interno de
        Control en el Instituto.
\end{enumerate}

El proceso de recepción de quejas descrito en el Artículo 13 establece que el
personal de la Defensoría debe determinar en primera instancia si la queja es
competente o incompetente. Esta decisión, aparentemente sencilla, presenta en la
práctica alta ambigüedad derivada de la redacción de los hechos. Dicha
dependencia del criterio manual genera tres conflictos críticos:

\begin{enumerate}
  \item El consumo excesivo de tiempo administrativo.
  \item La variabilidad de interpretaciones ante casos similares.
  \item La saturación operativa debida al incremento constante de solicitudes.
\end{enumerate}

En consecuencia, el problema se define como una tarea de \textit{Procesamiento
de Lenguaje Natural} (PLN) y \textit{clasificación supervisada}, donde el sistema
debe aprender a identificar los patrones lingüísticos y normativos que separan
una queja competente de una incompetente.

% ·············································································
\subsection{Taxonomía de denuncias}
% ·············································································

Las denuncias recibidas por la DDP no se clasifican únicamente por su contenido
textual, sino por una estructura jerárquica de validación normativa que el
sistema debe aprender a reconocer; por ello, establecer una taxonomía clara es
condición previa para construir cualquier modelo de clasificación. Transformar
una narrativa jurídica en un conjunto de datos procesable por algoritmos de
aprendizaje automático requiere organizar el conocimiento normativo de forma
sistemática. En el contexto de la Defensoría, esta taxonomía se apoya en dos
reglas críticas:

\begin{enumerate}
  \item \textbf{Pertenencia a la comunidad y jerarquía.} Define si las personas
        involucradas pertenecen a la comunidad politécnica. El sistema deberá
        identificar la existencia de un quejoso con su estatus institucional y
        la autoridad a quien se imputa la violación. Sin esta validación de
        identidad, la queja carece de base para ser procesada.

  \item \textbf{Exclusiones del Artículo 4°.} Constituye el núcleo del
        clasificador. Las denuncias se categorizan según la presencia de los
        cinco criterios de exclusión mencionados anteriormente. Esta regla se
        traduce en el sistema como siete campos binarios —características
        estructuradas— que procesan el análisis normativo:
        \begin{itemize}
          \item Naturaleza laboral vs. derechos politécnicos.
          \item Evaluación académica pura vs. vulneración de derechos en
                evaluación.
          \item Afectación individual vs. colectiva.
          \item Resolución disciplinaria vs. acto de autoridad.
          \item Procedimientos del OIC vs. acto u omisión de autoridad
                institucional.
        \end{itemize}
\end{enumerate}

Esta estructuración permite que el problema de clasificación no dependa
únicamente del texto libre, sino de reglas establecidas en el Artículo 4°. Al
final del proceso taxonómico, el sistema asigna a cada expediente una etiqueta
binaria definitiva: \textit{Competente} o \textit{No Competente} (canalización
a otra instancia o desechamiento), lo que constituye la variable de salida de los
modelos de clasificación supervisada.

% ·············································································
\subsection{Clasificación binaria y lógica de admisión}
% ·············································································

Una vez definida la taxonomía normativa, el siguiente paso es determinar la
forma en que el clasificador producirá su veredicto; este sistema adopta un
esquema de \textbf{clasificación binaria}, decisión fundamentada directamente en
el Artículo 13 del Acuerdo número extraordinario 622. A diferencia de otros
sistemas de categorización temática —que discriminan entre dominios como
deportes, política o finanzas— el clasificador de la DDP produce un veredicto
dicotómico: \textit{Competente} o \textit{No Competente}. Esta simplificación a
dos clases responde a la lógica de admisión procesal por tres razones
fundamentales:

\begin{enumerate}
  \item \textbf{Naturaleza del flujo administrativo.} Dentro de la DDP, el primer
        contacto actúa como una compuerta. No existe un estado intermedio: una
        queja o es admitida para su investigación, o es desechada y canalizada a
        otra instancia. La clasificación binaria representa fielmente este proceso
        de toma de decisiones institucional.

  \item \textbf{Optimización del entrenamiento.} Desde la perspectiva del
        aprendizaje automático, concentrar la muestra en dos etiquetas permite
        una mayor densidad de datos por categoría. Dado que el volumen de
        expedientes históricos etiquetados es finito, un modelo binario es más
        robusto y menos propenso al error que un modelo multiclase que intentara
        identificar simultáneamente el tipo específico de exclusión: laboral,
        académica o disciplinaria.

  \item \textbf{Reducción del falso negativo.} La lógica de admisión prioriza
        la protección de los derechos de la comunidad. Al configurar el problema
        como binario, es posible ajustar el umbral de decisión del modelo
        —especialmente en algoritmos como SVM— para minimizar los falsos
        negativos, asegurando que ninguna queja potencialmente competente sea
        descartada por una clasificación errónea.
\end{enumerate}

Por lo tanto, si el modelo predice la clase \textit{No Competente}, el sistema
asiste al personal sugiriendo la fundamentación legal basada en el Artículo 4°.
Si la predicción es \textit{Competente}, el expediente fluye automáticamente hacia
la etapa de radicación en la Subdefensoría, logrando así la meta de reducir los
tiempos de respuesta institucional.

% ·············································································
\subsection{Ciclo de vida de los datos en el aprendizaje supervisado}
% ·············································································

Para que la clasificación binaria opere de manera confiable, la narrativa jurídica
del quejoso debe recorrer un flujo técnico que la transforme en una señal numérica
apta para el aprendizaje automático; en consecuencia, el ciclo de vida de los
datos estructura en seis etapas el puente entre el texto libre y la decisión del
modelo. Este ciclo asegura la trazabilidad y la integridad de la información
jurídica en cada una de sus fases:

\begin{enumerate}
  \item \textbf{Captura de datos.} El ciclo inicia en la interfaz de registro
        donde el quejoso proporciona de forma estructurada su identidad
        institucional y, de forma no estructurada, el relato de los hechos.
        Esta etapa asegura que la entrada cumpla con los requisitos mínimos de
        admisión definidos en el Artículo 13.

  \item \textbf{Preprocesamiento y limpieza.} La narrativa libre se somete a
        técnicas de PLN. Se eliminan elementos de ruido, se normaliza el texto
        y se identifican términos clave que corresponden a las exclusiones
        normativas del Artículo 4°.

  \item \textbf{Vectorización y enriquecimiento.} El texto se convierte en
        representaciones vectoriales (\textit{embeddings}). Simultáneamente,
        el sistema inyecta las características jurídicas binarias extraídas del
        expediente, aplicando el factor de amplificación $\alpha$ para dar peso
        a la señal normativa sobre la semántica.

  \item \textbf{Entrenamiento y ajuste.} Utilizando un conjunto de datos
        histórico previamente etiquetado, el modelo aprende a identificar las
        fronteras de decisión legal. Se emplean técnicas de validación cruzada
        para medir la robustez de la generalización.

  \item \textbf{Inferencia y sugerencia.} Ante una queja nueva, el modelo
        procesa los vectores y genera una predicción de competencia acompañada
        de un nivel de probabilidad. El sistema asiste al personal humano
        mostrando el veredicto sugerido y la posible fundamentación legal.

  \item \textbf{Retroalimentación.} El ciclo se cierra cuando el personal
        jurídico valida o corrige la sugerencia del modelo. Estas decisiones
        se reincorporan al repositorio de datos para el reentrenamiento futuro,
        permitiendo que el clasificador evolucione conforme cambian los criterios
        institucionales.
\end{enumerate}

% ─────────────────────────────────────────────────────────────────────────────
\section{Fundamentos de Procesamiento de Lenguaje Natural}
% ─────────────────────────────────────────────────────────────────────────────

El \textbf{Procesamiento de Lenguaje Natural} (PLN) es la rama de la
inteligencia artificial que estudia cómo dotar a las computadoras de la capacidad
de leer, interpretar y generar texto humano. Su importancia radica en que el
lenguaje natural —a diferencia de los datos numéricos o tabulares— es inherentemente
ambiguo, dependiente del contexto y rico en matices que una máquina no puede
procesar directamente. \cite{jurafsky2023} Para que un algoritmo pueda razonar sobre un texto, este
debe primero transformarse en una representación matemática que preserve su
significado.

En el contexto de este proyecto, el PLN cumple un rol central: las quejas que
recibe la DDP llegan como narrativas libres, redactadas en lenguaje coloquial o
semijurídico por personas sin formación legal. Esto introduce desafíos
particulares que no están presentes en dominios más estandarizados: el
vocabulario especializado varía entre quejosos, la estructura argumentativa puede
ser implícita, y la clasificación jurídica correcta depende no solo de las
palabras usadas sino de la relación normativa que existe entre los hechos
descritos y los artículos del Acuerdo 622. Por lo tanto, dos quejas léxicamente
similares pueden tener clasificaciones opuestas —una competente y otra no— según
el contexto institucional que las rodea.

Para abordar estos desafíos, el sistema requiere dos tipos de herramientas PLN
que se desarrollan en las subsecciones siguientes: primero, el fundamento
matemático que permite representar documentos como vectores comparables; y
después, las transformaciones lingüísticas que preparan y limpian el texto antes
de que ese fundamento matemático pueda aplicarse.

% ·············································································
\subsection{Álgebra lineal en espacios vectoriales de alta dimensión}
% ·············································································

Para que una computadora pueda comparar documentos, determinar cuáles son
similares entre sí o trazar una frontera que separe categorías, primero necesita
que esos documentos existan como objetos matemáticos. El \textbf{álgebra lineal}
provee ese lenguaje: permite representar cualquier texto como un vector —una
lista ordenada de números— y luego medir relaciones entre vectores con precisión
aritmética. Esta representación es el fundamento sobre el que operan todos los
vectorizadores y clasificadores empleados en este sistema; sin ella, no existiría
un mecanismo para que el modelo aprenda qué quejas son parecidas entre sí ni cómo
trazar la línea que separa las competentes de las incompetentes.

La idea central es que cualquier documento puede representarse como un punto en
un \textbf{espacio vectorial de alta dimensión} $\mathbb{R}^d$, donde $d$ es el
número de características que describen al documento. Documentos semánticamente
similares —por ejemplo, dos quejas sobre acoso académico— ocuparán regiones
cercanas en ese espacio; documentos disímiles —una queja laboral y una queja por
evaluación— quedarán alejados. La \textbf{similitud coseno} formaliza esta noción
de cercanía entre dos documentos $\mathbf{u}$ y $\mathbf{v}$:

\begin{equation}
  \text{sim}(\mathbf{u}, \mathbf{v})
  = \frac{\mathbf{u} \cdot \mathbf{v}}{\|\mathbf{u}\| \, \|\mathbf{v}\|}
  = \frac{\displaystyle\sum_{i=1}^{d} u_i v_i}
         {\sqrt{\displaystyle\sum_{i=1}^{d} u_i^2}
          \cdot \sqrt{\displaystyle\sum_{i=1}^{d} v_i^2}}
  \label{eq:coseno}
\end{equation}

donde $\text{sim} \in [-1, 1]$. Un valor cercano a $1$ indica alta similitud
semántica; esta métrica es la base del razonamiento por proximidad que subyace
a todos los vectorizadores empleados en el sistema.

La elección de cuántas dimensiones usar para representar los documentos no es
arbitraria. En este proyecto se trabaja con:

\begin{itemize}
\item 300 dimensiones esto es por SpaCy la cual fue entrenada por Explosion AI usando word2vec y esta última fue entrenada originalmente por Mikolov et al. (2013) con vectores de 300 dimensiones, conrvitiendose en el estándar de la industria.
\item 768 dimensiones por SentenceBert. Ese número viene de la arquitectura BERT base, que usa capas de atención con una dimensión oculta de 768. Cuando Reimers y Gurevych construyeron SBERT sobre BERT, el vector de salida del pooling final hereda esa dimensión.
\end{itemize}

A mayor dimensión, el modelo puede
capturar matices semánticos más finos —por ejemplo, distinguir entre
\textit{``no se me permitió acceder''} y \textit{``accedí sin autorización''}—
pero también requiere mayor capacidad computacional y datos suficientes para
evitar la \textit{maldición de la dimensionalidad}: conforme crece $d$, el
volumen del espacio aumenta exponencialmente y los datos se vuelven muy dispersos,
lo que dificulta encontrar vecinos significativos. Sobre estos espacios, el
sistema ejecuta cuatro operaciones elementales del álgebra lineal:

\begin{itemize}
  \item \textbf{Suma de vectores:} permite combinar representaciones de tokens
        individuales en una representación de documento.
  \item \textbf{Promedio ponderado:} base del vectorizador de spaCy.
  \item \textbf{Producto punto y norma:} usados internamente por SVM para
        calcular márgenes de separación.
  \item \textbf{Concatenación horizontal:} operación central de la arquitectura
        híbrida para unir embeddings de texto con features jurídicos.
\end{itemize}

Sin embargo, para que estas operaciones produzcan vectores con significado real,
el texto de entrada debe estar limpio y estandarizado. Un modelo no puede extraer
un vector confiable de una queja llena de errores tipográficos, abreviaturas
inconsistentes o términos sinónimos tratados como palabras distintas. De ahí la
necesidad del preprocesamiento descrito en la subsección siguiente.

% ·············································································
\subsection{Preprocesamiento de la queja}
% ·············································································

Antes de que el álgebra lineal pueda operar sobre el texto, la narrativa libre
del quejoso debe pasar por una cadena de transformaciones que reduzcan el ruido
y estandaricen la representación lingüística; este preprocesamiento es
especialmente relevante en textos jurídicos, donde variaciones ortográficas,
errores tipográficos o sinónimos coloquiales de términos normativos pueden afectar
directamente la calidad de los embeddings. Las transformaciones se aplican en el
siguiente orden:

\subsubsection*{Tokenización}

La \textbf{tokenización} segmenta el texto continuo en unidades mínimas de
significado (tokens) \cite{jurafsky2023}. SpaCy realiza una tokenización sensible al contexto que
distingue, por ejemplo, contracciones y signos de puntuación:

\begin{center}
  \texttt{"me negó el acceso"} $\longrightarrow$
  [\texttt{me}, \texttt{negó}, \texttt{el}, \texttt{acceso}]
\end{center}

\subsubsection*{Lematización}

La \textbf{lematización} reduce cada token a su forma canónica (\textit{lema}),
eliminando variaciones morfológicas \cite{jurafsky2023}:

\begin{center}
  \texttt{negaron} $\to$ \texttt{negar}, \quad
  \texttt{discriminado} $\to$ \texttt{discriminar}, \quad
  \texttt{estudiantiles} $\to$ \texttt{estudiantil}
\end{center}

\subsubsection*{Filtrado de palabras vacías}

Las \textbf{palabras vacías} (\textit{stop words}) son términos de alta
frecuencia y bajo contenido semántico (artículos, preposiciones, conjunciones).
En textos jurídicos, sin embargo, ciertas palabras que en otros contextos serían
vacías — como \textit{``no''}, \textit{``sin''} o \textit{``nunca''} — pueden
ser jurídicamente decisivas. spaCy permite personalizar esta lista para el
dominio específico de la Defensoría.

\subsubsection*{Análisis morfosintáctico con spaCy}

El modelo \texttt{es\_core\_news\_lg} de spaCy realiza, además de la
tokenización y lematización, un análisis morfosintáctico completo que asigna a
cada token su categoría gramatical (\textit{POS tagging}) y sus dependencias
sintácticas. Esto permite identificar, por ejemplo, que en la oración:

\begin{center}
  \textit{``El director me negó el acceso al laboratorio''}
\end{center}

\noindent el sujeto es \textit{director} (autoridad), el verbo es \textit{negar}
(acto de omisión) y el complemento es \textit{acceso al laboratorio} (derecho
afectado). Esta información estructural enriquece la representación semántica
más allá del léxico superficial, y sienta las bases para la etapa de
vectorización descrita en la sección siguiente.

% ─────────────────────────────────────────────────────────────────────────────
\section{Representación Vectorial del Texto}
% ─────────────────────────────────────────────────────────────────────────────

Una vez que el texto ha sido normalizado mediante el preprocesamiento, el siguiente
desafío consiste en convertirlo en representaciones numéricas que los algoritmos
de aprendizaje automático puedan procesar; la representación vectorial del texto
es, por lo tanto, el puente entre el lenguaje natural y los clasificadores. A
continuación se presenta la evolución de estas técnicas desde los modelos
estadísticos hasta los modelos de lenguaje contextuales de última generación,
mostrando en cada etapa por qué los enfoques más sencillos son insuficientes para
el dominio jurídico de la DDP.

% ·············································································
\subsection{Modelos de Bolsa de Palabras y TF-IDF}
% ·············································································

Los modelos estadísticos de representación textual constituyen el punto de
partida histórico de la vectorización y permiten comprender, por contraste, las
limitaciones que motivan el uso de embeddings semánticos en este sistema. Los
\textbf{modelos de Bolsa de Palabras} (\textit{Bag of Words}, BoW) representan
un documento como un vector de frecuencias de términos sobre el vocabulario total
del corpus. En BoW cada palabra individual se convierte en una dimensión separada
del espacio vectorial \cite{manning2008}: si un conjunto de texto tiene $n$ palabras, el espacio
vectorial resultante tendrá $n$ dimensiones, y la posición de un documento en ese
espacio está determinada por la cantidad de veces que aparece cada término.
Formalmente:

\begin{equation}
  \mathbf{v}_{BoW}(d) = [f(t_1, d),\; f(t_2, d),\; \ldots,\; f(t_V, d)]
  \label{eq:bow}
\end{equation}

donde $f(t_i, d)$ es la frecuencia del término $t_i$ en el documento $d$ y $V$
es el tamaño del vocabulario.

\textbf{TF-IDF} (\textit{Term Frequency -- Inverse Document Frequency}) refina
este enfoque ponderando cada término según su frecuencia en el documento y su
rareza en el corpus \cite{manning2008}:

\begin{equation}
  \text{tf-idf}(t, d, D)
  = \underbrace{\frac{f(t, d)}{\sum_{t'} f(t', d)}}_{\text{TF}}
    \times
    \underbrace{\log \frac{|D|}{|\{d' \in D : t \in d'\}|}}_{\text{IDF}}
  \label{eq:tfidf}
\end{equation}

\noindent\textbf{Limitaciones críticas para el dominio jurídico.} Aunque son
computacionalmente eficientes, estos modelos son inadecuados para clasificar
quejas de la DDP por dos razones fundamentales:

\begin{enumerate}
  \item \textbf{Pérdida del orden y contexto:} Las frases \textit{``el director
        negó el acceso''} y \textit{``el acceso negó al director''} generan el
        mismo vector, aunque tienen significados radicalmente distintos.
  \item \textbf{Sesgo léxico superficial:} El término \textit{``calificación''}
        tendría alta frecuencia tanto en quejas por evaluación académica (No
        Competente) como en quejas por extorsión académica (Competente), haciendo
        imposible la distinción correcta.
\end{enumerate}

Estos modelos se presentan como antecedentes necesarios para contextualizar la
superioridad de los enfoques basados en embeddings empleados en este sistema.

% ·············································································
\subsection{Word Embeddings Estáticos — spaCy (\texttt{es\_core\_news\_lg})}
% ·············································································

Dado que los modelos estadísticos pierden el contexto semántico, la siguiente
generación de técnicas representa cada palabra como un vector denso aprendido
de grandes corpus; esta representación, conocida como \textbf{embedding
estático}, resuelve la limitación semántica de BoW al colocar palabras
relacionadas —como \textit{acoso} y \textit{hostigamiento}— en regiones cercanas
del espacio vectorial. Los embeddings estáticos representan cada palabra como un
vector denso en $\mathbb{R}^{300}$, entrenado sobre grandes corpus en español
mediante el algoritmo \textit{word2vec}.\cite{mikolov2013}

El modelo \texttt{es\_core\_news\_lg} de spaCy \cite{honnibal2020} fue entrenado sobre el corpus
\textit{Spanish Web Crawl} y genera vectores de 300 dimensiones por token. La
representación de un documento completo se obtiene promediando los vectores de
los tokens informativos:

\begin{equation}
  \mathbf{v}_{spaCy}(d)
  = \frac{1}{|T|} \sum_{t \in T} \mathbf{w}_t, \qquad
  T = \bigl\{t \in d \;\big|\; \texttt{has\_vector}(t)
      \wedge \neg\,\texttt{is\_stop}(t)\bigr\}
  \label{eq:spacy}
\end{equation}

donde $\mathbf{w}_t \in \mathbb{R}^{300}$ es el vector del token $t$ y $T$ es
el conjunto de tokens informativos (con vector disponible y no vacías).

No obstante, al promediar todos los vectores del documento se pierde el orden
sintáctico y el contexto local. La frase \textit{``el profesor me acosó
condicionando la calificación''} y la frase \textit{``el profesor condicionó la
calificación — sin acosar''} generarían vectores cercanos, a pesar de tener
clasificaciones jurídicas distintas. Esta limitación motiva el uso de modelos
contextuales como Sentence-BERT, presentados a continuación.

% ·············································································
\subsection{Sentence Transformers — SBERT paraphrase-multilingual-mpnet-base-v2}
% ·············································································

Para superar la insensibilidad al orden y al contexto de los embeddings estáticos,
se requiere un modelo que genere una representación de toda la oración tomando en
cuenta las relaciones entre sus partes; \textbf{Sentence-BERT (SBERT)} satisface
este requisito y representa el estado del arte en codificación semántica de
textos.\cite{reimers2019} A diferencia de los embeddings estáticos, SBERT genera
un único vector de representación para toda la oración, capturando las relaciones
entre palabras a lo largo del texto completo mediante la arquitectura
\textit{Transformer}.

\subsubsection*{Arquitectura Transformer}

El bloque central del Transformer es el mecanismo de \textbf{atención
multi-cabeza} (\textit{multi-head self-attention}), que calcula para cada token
su relevancia en función de todos los demás tokens del documento:

\begin{equation}
  \text{Attention}(Q, K, V)
  = \text{softmax}\!\left(\frac{QK^\top}{\sqrt{d_k}}\right) V
  \label{eq:attention}
\end{equation}

donde $Q$, $K$ y $V$ son las matrices de consulta, clave y valor derivadas
de los embeddings de entrada, y $d_k$ es la dimensión de la clave.\cite{vaswani2017} Este
mecanismo permite que el modelo asigne mayor peso a los tokens contextualmente
relevantes. Por ejemplo, en la queja:

\begin{center}
  \textit{``Me reprobó porque no acepté salir con él''}
\end{center}

\noindent el mecanismo de atención aprende a conectar \textit{reprobó} con
\textit{acepté salir}, identificando que la relación causal entre ambos
constituye un acto de acoso, no una evaluación académica pura.

\subsubsection*{Entrenamiento contrastivo de SBERT}

SBERT extiende BERT con una capa de \textit{pooling} sobre los estados ocultos
y se entrena mediante aprendizaje contrastivo con pares de oraciones
semánticamente equivalentes y no equivalentes:

\begin{equation}
  \mathcal{L}_{contrastiva}
  = \sum_{(i,j)} \bigl[\text{sim}(\mathbf{e}_i, \mathbf{e}_j) - y_{ij}\bigr]^2
  \label{eq:contrastiva}
\end{equation}

donde $y_{ij} = 1$ si las oraciones $i$ y $j$ son semánticamente equivalentes
y $y_{ij} = 0$ en caso contrario.

\subsubsection*{Modelo empleado}

El modelo \texttt{paraphrase-multilingual-mpnet-base-v2} fue entrenado sobre
más de 50 idiomas, incluido el español. Genera embeddings de \textbf{768
dimensiones} por documento y ha demostrado superioridad frente a los promedios
de \textit{word embeddings} en tareas de clasificación de texto corto y mediano.
Esta capacidad es crítica para distinguir los casos límite de la DDP, donde
el acto jurídicamente relevante puede estar expresado en una sola oración dentro
de un párrafo más extenso.

% ─────────────────────────────────────────────────────────────────────────────
\section{Ingeniería de Características Híbridas}
% ─────────────────────────────────────────────────────────────────────────────

Aunque los embeddings semánticos capturan el significado del texto, los modelos
puramente textuales pueden aprender correlaciones léxicas espurias en lugar de la
lógica jurídica subyacente; por ello, una de las contribuciones centrales de este
sistema es la \textbf{arquitectura híbrida} que combina la representación
semántica del texto con conocimiento normativo estructurado. Esta arquitectura
opera en dos pasos: primero codifica los criterios del Artículo 4° como variables
binarias, y después las fusiona con el embedding textual mediante una
concatenación amplificada que equilibra ambas señales.

% ·············································································
\subsection{Codificación de variables normativas}
% ·············································································

Los cinco criterios de incompetencia del Artículo 4° se traducen en
\textbf{siete variables binarias} (features jurídicos) que codifican el análisis
normativo de cada queja de forma estructurada; esta codificación es el mecanismo
que permite al clasificador razonar con base en la ley y no solo en la semántica
del texto. Las variables y su descripción normativa se presentan en la Tabla
\ref{tab:features_mt}:

\begin{table}[H]
\centering
\caption{Features jurídicos estructurados derivados del Artículo 4°.}
\label{tab:features_mt}
\begin{tabularx}{\textwidth}{clX}
\toprule
\textbf{\#} & \textbf{Variable} & \textbf{Descripción normativa} \\
\midrule
1 & \texttt{es\_conflicto\_laboral\_art4\_I} &
    El fondo del conflicto es una relación laboral o sindical
    (salarios, plazas, escalafón, prestaciones). \\
2 & \texttt{es\_evaluacion\_academica\_pura\_art4\_II} &
    La queja versa \emph{únicamente} sobre la valoración del
    aprendizaje por parte de un profesor o comisión. \\
3 & \texttt{hay\_vulneracion\_notoria\_en\_evaluacion} &
    Existe una violación de derechos dentro de un proceso
    evaluativo (extorsión, acoso, discriminación). \\
4 & \texttt{es\_resolucion\_disciplinaria\_art4\_III} &
    La queja impugna una resolución disciplinaria ya emitida. \\
5 & \texttt{es\_afectacion\_colectiva\_art4\_IV} &
    La afectación descrita es grupal, no individual. \\
6 & \texttt{involucra\_oic\_art4\_V} &
    El asunto está siendo conocido por el Órgano Interno de
    Control. \\
7 & \texttt{hay\_acto\_u\_omision\_autoridad} &
    Existe un acto u omisión de una autoridad institucional
    que vulnera derechos politécnicos. \\
\bottomrule
\end{tabularx}
\end{table}

Estas variables se representan como un vector $\mathbf{f} \in \{0, 1\}^7$ y
se determinan a partir del análisis del expediente, ya sea de forma manual
durante la construcción del corpus o de forma automática en una etapa futura
mediante un modelo extractor. La lógica de consistencia entre las variables y
la clase asignada es la siguiente:

\begin{itemize}
  \item Si \texttt{clase = No Competente} por razón laboral:
        \texttt{es\_conflicto\_laboral\_art4\_I = 1}, resto $= 0$.
  \item Si \texttt{clase = Competente}: \texttt{hay\_acto\_u\_omision\_autoridad = 1},
        todas las fracciones del Art. 4° $= 0$.
  \item La variable 3 (\texttt{hay\_vulneracion\_notoria}) solo puede ser $1$
        cuando la variable 2 (\texttt{es\_evaluacion\_academica\_pura}) también
        es $1$ \emph{y} la clase es \textit{Competente}: esto modela la excepción
        de la fracción IV del Artículo 4°.
\end{itemize}

% ·············································································
\subsection{Concatenación y Factor de Amplificación ($\alpha$)}
% ·············································································

Una vez codificadas las variables normativas, es necesario fusionarlas con el
embedding semántico de manera que ambas señales contribuyan equitativamente a la
decisión del clasificador; para ello, la unión se realiza mediante
\textbf{concatenación horizontal} junto con un factor de amplificación $\alpha$
que compensa la diferencia de escala entre los vectores:

\begin{equation}
  \mathbf{x}_{final}
  = \left[\mathbf{e} \;\|\; \alpha \cdot \mathbf{f}\right]
  \in \mathbb{R}^{d + 7}
  \label{eq:concat}
\end{equation}

donde $d \in \{300, 768\}$ según el vectorizador empleado y $\alpha$ es el
\textbf{factor de amplificación}.

\subsubsection*{Justificación matemática del factor $\alpha$}

Sin amplificación ($\alpha = 1$), los 7 features jurídicos quedan efectivamente
opacados por las $d$ dimensiones del embedding, cuya norma promedio es
significativamente mayor. Para ilustrarlo, considérese la contribución relativa
de cada componente a la norma del vector concatenado:

\begin{equation}
  \frac{\|\alpha \cdot \mathbf{f}\|}{\|\mathbf{e}\| + \|\alpha \cdot \mathbf{f}\|}
  = \frac{\alpha \sqrt{7}}{\|\mathbf{e}\| + \alpha \sqrt{7}}
  \label{eq:contribucion}
\end{equation}

Con $\|\mathbf{e}\| \approx 15$ (valor típico para embeddings de Sentence-BERT
normalizados) y $\alpha = 1$:

\begin{equation*}
  \frac{1 \cdot \sqrt{7}}{15 + \sqrt{7}} \approx \frac{2.65}{17.65} \approx 15\%
\end{equation*}

Con $\alpha = 3$:

\begin{equation*}
  \frac{3\sqrt{7}}{15 + 3\sqrt{7}} \approx \frac{7.94}{22.94} \approx 35\%
\end{equation*}

El valor $\alpha = 3.0$ fue seleccionado experimentalmente como el punto de
equilibrio donde la señal jurídica tiene peso suficiente para influir en las
fronteras de decisión del clasificador, sin dominar completamente sobre la
información semántica del texto.

% ·············································································
\subsection{Simulación de producción: entrenamiento con metadatos
            e inferencia en modo autónomo}
% ·············································································

La arquitectura híbrida introduce una asimetría importante entre el escenario de
entrenamiento y el escenario de producción, que debe modelarse explícitamente
para evitar que el sistema sea sobreoptimista en sus métricas de evaluación; esta
asimetría surge porque, durante la construcción del corpus, el personal jurídico
anota manualmente los siete atributos booleanos del Artículo 4°, mientras que en
producción el sistema recibe únicamente el relato libre del quejoso sin anotación
previa. Esta asimetría se modela mediante dos modos de operación diferenciados:

\begin{enumerate}
  \item \textbf{Modo asistido.} El personal de primer contacto completa el
        formulario con los atributos jurídicos mientras atiende al quejoso. El
        clasificador opera con el vector completo
        $\mathbf{x}_{final} = [\mathbf{e} \| \alpha \cdot \mathbf{f}]$,
        aprovechando tanto la semántica del texto como la señal normativa
        estructurada.

  \item \textbf{Modo autónomo.} Los atributos jurídicos se fijan en cero
        ($\mathbf{f} = \mathbf{0}$), de forma que el clasificador basa su
        veredicto exclusivamente en la representación semántica del texto:
        $\mathbf{x}_{final} = [\mathbf{e} \| \mathbf{0}]$.
        Este modo representa el escenario de producción real y constituye la
        condición de estrés bajo la cual se evalúa la robustez de los modelos.
\end{enumerate}

La distinción entre ambos modos tiene implicaciones directas para la selección
del clasificador óptimo: un modelo que dependa excesivamente de la señal
normativa amplificada puede exhibir métricas perfectas durante el entrenamiento,
pero degradarse significativamente en producción cuando dicha señal no está
disponible. Por esta razón, el criterio de selección del modelo final privilegia
el desempeño en modo autónomo sobre el desempeño con metadatos.

% ─────────────────────────────────────────────────────────────────────────────
\section{Modelos de Aprendizaje Automático Evaluados}
% ─────────────────────────────────────────────────────────────────────────────

Con el vector de características híbridas definido, el siguiente paso es elegir
el algoritmo que mejor aprenda las fronteras de decisión legal a partir de él;
para garantizar la selección del clasificador óptimo para el dominio jurídico de
la DDP, se evaluaron cuatro configuraciones de algoritmos de aprendizaje
supervisado con distintos paradigmas de construcción de fronteras de decisión.
Cada modelo aporta una perspectiva diferente sobre los datos, y su comparación
permite identificar cuál se adapta mejor a la distribución de las quejas en el
espacio vectorial.

% ·············································································
\subsection{Máquinas de Soporte Vectorial (SVM)}
% ·············································································

Las \textbf{Máquinas de Soporte Vectorial} (\textit{Support Vector Machines},
SVM) son algoritmos de clasificación que buscan el \textbf{hiperplano de máximo
margen} que separa las clases en el espacio de características, lo que las hace
especialmente robustas ante casos límite como las quejas que se ubican en la
frontera entre competente e incompetente. \cite{cortes1995} Dado un conjunto de
entrenamiento $\{(\mathbf{x}_i, y_i)\}_{i=1}^n$ con $y_i \in \{-1, +1\}$, la
SVM resuelve el problema de optimización:

\begin{equation}
  \min_{\mathbf{w}, b, \boldsymbol{\xi}}
  \frac{1}{2}\|\mathbf{w}\|^2 + C \sum_{i=1}^n \xi_i
  \quad \text{s.a.} \quad
  y_i(\mathbf{w}^\top \mathbf{x}_i + b) \geq 1 - \xi_i, \quad \xi_i \geq 0
  \label{eq:svm}
\end{equation}

donde $C$ es el parámetro de regularización que controla la tolerancia al error
de clasificación ($C = 0.5$ en este sistema) y $\boldsymbol{\xi}$ son variables
de holgura que permiten clasificación blanda.

\subsubsection*{Kernel RBF}

Para datos no linealmente separables en el espacio original, la SVM proyecta
los datos a un espacio de mayor dimensión mediante la \textbf{función de kernel}.
El kernel de Base Radial (\textit{Radial Basis Function}, RBF) empleado en este
sistema calcula la similitud entre dos puntos como:

\begin{equation}
  K(\mathbf{x}_i, \mathbf{x}_j)
  = \exp\!\left(-\gamma \|\mathbf{x}_i - \mathbf{x}_j\|^2\right)
  \label{eq:rbf}
\end{equation}

donde $\gamma$ controla el radio de influencia de cada vector de soporte. El
kernel RBF es especialmente adecuado para embeddings de alta dimensión porque
puede construir fronteras de decisión no lineales arbitrariamente complejas,
adaptándose a la distribución de los casos jurídicos en el espacio semántico \cite{cortes1995}.

\subsubsection*{Ventaja en el dominio jurídico}

La SVM con kernel RBF maximiza el margen de separación entre las clases, lo
que le confiere una robustez natural ante casos límite: quejas que caen en la
frontera entre \textit{Competente} y \textit{No Competente} se clasifican
según el hiperplano de máximo margen, no según la densidad local de puntos.
No obstante, esta ventaja geométrica está condicionada a la disponibilidad de
una señal de entrada suficientemente discriminativa; cuando los atributos
jurídicos no están presentes, la eficacia del clasificador depende enteramente
de la calidad del embedding semántico.

% ·············································································
\subsection{Regresión Logística}
% ·············································································

A diferencia de la SVM —que construye una frontera geométrica de máximo margen—
la \textbf{Regresión Logística} es un modelo lineal que estima directamente la
probabilidad de pertenencia a cada clase, lo que le confiere una interpretabilidad
valiosa en el contexto jurídico donde la explicabilidad de la decisión puede ser
requerida institucionalmente. La probabilidad de pertenencia a la clase positiva
se estima mediante la función logística (sigmoide):

\begin{equation}
  P(y = 1 \mid \mathbf{x})
  = \sigma(\mathbf{w}^\top \mathbf{x} + b)
  = \frac{1}{1 + e^{-(\mathbf{w}^\top \mathbf{x} + b)}}
  \label{eq:logistica}
\end{equation}

Los parámetros $\mathbf{w}$ y $b$ se estiman minimizando la función de pérdida
de entropía cruzada binaria:

\begin{equation}
  \mathcal{L}(\mathbf{w}, b)
  = -\frac{1}{n} \sum_{i=1}^n
    \Bigl[y_i \log \hat{p}_i + (1 - y_i) \log(1 - \hat{p}_i)\Bigr]
  + \lambda \|\mathbf{w}\|^2
  \label{eq:cross-entropy}
\end{equation}

donde $\lambda$ es el término de regularización L2. La configuración empleada
en este sistema es \texttt{max\_iter=1000} con $C = 0.3$. Adicionalmente, los
coeficientes $w_j$ son directamente interpretables: un coeficiente positivo
asociado a una dimensión del embedding indica que ese componente favorece la
clasificación como \textit{Competente}, lo que facilita la auditoría del modelo.

% ·············································································
\subsection{Ensamble de Árboles — Random Forest}
% ·············································································

Donde la SVM y la Regresión Logística construyen una sola frontera de decisión,
el \textbf{Random Forest} resuelve el problema mediante la agregación de múltiples
árboles de decisión entrenados sobre subconjuntos aleatorios de los datos; esta
estrategia de ensamble reduce la varianza de la predicción individual y confiere
al modelo una notable capacidad de generalización. \cite{breiman2001} Formalmente:

\begin{equation}
  \hat{y}_{RF}(\mathbf{x})
  = \frac{1}{B} \sum_{b=1}^{B} T_b(\mathbf{x})
  \label{eq:rf}
\end{equation}

donde $B$ es el número de árboles ($B = 300$ en este sistema) y $T_b$ es el
$b$-ésimo árbol entrenado sobre una muestra bootstrap. Los árboles de decisión
particionan el espacio de características con cortes ortogonales a los ejes,
construyendo regiones de decisión rectangulares de forma recursiva. Aunque en
espacios de dimensión muy alta esta estrategia puede ser menos eficiente que el
hiperplano de máximo margen de la SVM, el promedio de un número grande de árboles
con submuestreo aleatorio de características reduce significativamente la varianza
y resulta especialmente robusto cuando la representación semántica del texto no
está apoyada por atributos estructurados.

% ·············································································
\subsection{Ensemble por Votación Suave}
% ·············································································

Dado que cada clasificador individual tiene fortalezas complementarias, una
estrategia de \textit{ensemble} que combine sus predicciones puede alcanzar un
desempeño superior al de cualquiera de ellos por separado; la \textbf{votación
suave} (\textit{soft voting}) aprovecha esta propiedad promediando las
probabilidades predichas por múltiples clasificadores base para obtener una
predicción final más estable. \cite{dietterich2000}

\begin{equation}
  \hat{y} = \arg\max_{c} \frac{1}{M} \sum_{m=1}^{M} p_m(c \mid \mathbf{x})
  \label{eq:soft_voting}
\end{equation}

donde $M$ es el número de modelos base y $p_m(c \mid \mathbf{x})$ es la
probabilidad que el modelo $m$ asigna a la clase $c$. A diferencia de la
\textbf{votación dura} —que cuenta votos de clase— la votación suave aprovecha
la magnitud de la confianza de cada clasificador, lo que generalmente produce
predicciones más calibradas cuando los clasificadores base tienen tasas de error
no correlacionadas.

En este sistema, el \textit{ensemble} integra SVM-RBF, Regresión Logística y
Random Forest. La motivación es que cada clasificador tiene fortalezas
complementarias: SVM maximiza el margen geométrico, la Regresión Logística
produce probabilidades linealmente calibradas, y Random Forest captura fronteras
no lineales por partición del espacio. Si sus errores no están correlacionados,
el \textit{ensemble} reduce la varianza de la predicción individual respecto a
cualquiera de los modelos base. Para que el \texttt{VotingClassifier} pueda
operar en modo \texttt{voting="soft"}, todos los estimadores deben exponer el
método \texttt{predict\_proba}; dado que \texttt{SVC} no lo implementa de forma
nativa, se envuelve en \texttt{CalibratedClassifierCV}, que estima probabilidades
mediante validación cruzada interna sobre el conjunto de entrenamiento.

% ─────────────────────────────────────────────────────────────────────────────
\section{Estrategias de Entrenamiento y Validación Robusta}
% ─────────────────────────────────────────────────────────────────────────────

La selección del mejor clasificador no puede basarse en métricas obtenidas sobre
los mismos datos usados para entrenarlo, ya que esto produciría una estimación
optimista del desempeño real; por ello, las estrategias de entrenamiento y
validación robusta garantizan que las métricas reportadas reflejen la capacidad
de generalización real del sistema y no sean artefactos de una partición
particular de los datos.

% ·············································································
\subsection{Validación Cruzada $K$-Fold Estratificada}
% ·············································································

La \textbf{validación cruzada $K$-fold} es la estrategia central para estimar
el desempeño de los modelos de manera confiable: divide el conjunto de datos en
$K$ particiones de tamaño aproximadamente igual y rota el conjunto de evaluación
de manera que todos los datos sean usados para prueba exactamente una vez. En
cada iteración, $K-1$ particiones se usan para entrenamiento y la restante para
evaluación, obteniendo una estimación promedio del desempeño:

\begin{equation}
  \overline{F1}_{CV}
  = \frac{1}{K} \sum_{k=1}^{K} F1_k
  \label{eq:cv}
\end{equation}

La variante \textbf{estratificada} garantiza que la proporción de clases se
mantiene en cada pliegue, lo cual es crítico ante corpus con posible desbalance
entre quejas competentes e incompetentes. En este proyecto se empleó $K = 5$,
lo que significa que en cada iteración el 80\% de los datos se usa para
entrenamiento y el 20\% restante para evaluación. Además del promedio, la
desviación estándar entre pliegues ($\sigma_{CV}$) se usa como medida de
\textbf{estabilidad} del modelo:

\begin{equation}
  \sigma_{CV} = \sqrt{\frac{1}{K} \sum_{k=1}^{K} (F1_k - \overline{F1}_{CV})^2}
  \label{eq:sigma}
\end{equation}

Un $\sigma_{CV}$ cercano a cero indica que el modelo no es sensible a la
partición específica de los datos, condición deseable para un sistema de apoyo
en producción.

\noindent\textbf{Nota de implementación.} En la implementación de este proyecto,
la validación cruzada $K = 5$ se aplica exclusivamente sobre el conjunto de
entrenamiento (70\% del corpus, 367 casos), no sobre el dataset completo. Esto
significa que cada \textit{fold} evalúa aproximadamente 73 casos internos del
entrenamiento. Esta validación se complementa con un conjunto \textit{holdout}
independiente del 15\% (79 casos), completamente aislado durante el entrenamiento
y la selección de hiperparámetros, cuyo F1 ponderado —evaluado en modo autónomo,
sin atributos jurídicos— constituye el criterio definitivo de selección del
modelo.

% ·············································································
\subsection{Evaluación mediante Métricas de Inteligencia}
% ·············································································

Las métricas de evaluación deben reflejar las prioridades institucionales de la
DDP y no únicamente la exactitud global del modelo; en consecuencia, se emplean
métricas que distinguen entre los distintos tipos de error, asignando mayor peso
al error con mayor costo institucional: clasificar como incompetente una queja
que sí debía ser atendida. Para comprender estas métricas, es necesario
establecer primero los cuatro resultados posibles de una clasificación binaria,
organizados en la \textbf{matriz de confusión}:

\begin{itemize}
  \item \textbf{Verdadero Positivo (TP):} la queja es competente y el modelo
        la clasifica como competente.
  \item \textbf{Verdadero Negativo (TN):} la queja es incompetente y el modelo
        la clasifica como incompetente.
  \item \textbf{Falso Positivo (FP):} la queja es incompetente pero el modelo
        la clasifica como competente.
  \item \textbf{Falso Negativo (FN):} la queja es competente pero el modelo
        la clasifica como incompetente. En el contexto de la DDP, este es el
        error más grave: implica descartar una queja que debía ser atendida.
\end{itemize}

\subsubsection*{Accuracy}

La \textbf{exactitud} (\textit{accuracy}) mide la proporción de predicciones
correctas sobre el total de instancias:

\begin{equation}
  \text{Accuracy} = \frac{TP + TN}{TP + TN + FP + FN}
  \label{eq:accuracy}
\end{equation}

En el contexto de la DDP, donde las consecuencias de un falso negativo son más
graves que las de un falso positivo, la accuracy por sí sola no es suficiente.
Por ello, se complementa con las métricas de Precisión, Recall y F1-Score.

\subsubsection*{Precisión, Recall y F1-Score}

\begin{align}
  \text{Precisión} &= \frac{TP}{TP + FP}
  \label{eq:precision}\\[4pt]
\intertext{La \textbf{Precisión} mide qué proporción de las quejas clasificadas
como competentes realmente lo son; penaliza las falsas alarmas.}
  \text{Recall} &= \frac{TP}{TP + FN}
  \label{eq:recall}\\[4pt]
\intertext{El \textbf{Recall} mide qué proporción de las quejas realmente
competentes fueron identificadas correctamente; penaliza los casos omitidos.
En el contexto de la DDP el Recall tiene mayor peso institucional: un Recall
bajo implica que quejas legítimas están siendo descartadas sin revisión.}
  F1 &= \frac{2 \cdot \text{Precisión} \cdot \text{Recall}}
             {\text{Precisión} + \text{Recall}}
  \label{eq:f1}
\end{align}

El \textbf{F1-Score} es la media armónica entre Precisión y Recall, ofreciendo
un balance que penaliza los extremos donde una métrica es alta a costa de la
otra. El \textbf{F1 ponderado} (\textit{weighted F1}) extiende este cálculo a
ambas clases, ponderando por el número de instancias de cada una:

\begin{equation}
  F1_w = \sum_{c \in C} \frac{n_c}{N} \cdot F1_c
  \label{eq:f1w}
\end{equation}

donde $n_c$ es el número de instancias de la clase $c$ y $N$ el total. Esta
métrica es la principal para la selección del modelo óptimo en este sistema.

\subsubsection*{Análisis de la varianza $\sigma$}

La desviación estándar del F1 en validación cruzada ($\sigma_{CV}$) complementa
el valor medio al revelar la estabilidad del modelo. Un modelo con
$\overline{F1}_{CV} = 0.98$ y $\sigma_{CV} = 0.04$ puede ser menos confiable
que uno con $\overline{F1}_{CV} = 0.97$ y $\sigma_{CV} = 0.00$, pues el
primero es sensible a la distribución particular de los datos en cada pliegue.

% ·············································································
\subsection{Carga Perezosa y Serialización con \texttt{joblib}}
% ·············································································

Una vez seleccionado el modelo óptimo, este debe ser almacenado y puesto en
producción de manera que garantice la reproducibilidad de sus predicciones y
minimice el tiempo de arranque del sistema; para ello se combinan la
\textbf{carga perezosa} y la \textbf{serialización con \texttt{joblib}}, dos
patrones de diseño que optimizan la eficiencia operativa del clasificador en
producción.

La \textbf{carga perezosa} (\textit{lazy loading}) inicializa los recursos
computacionalmente costosos —como la carga del modelo de Sentence-BERT o el
modelo de spaCy— únicamente en el momento en que son requeridos por primera vez,
no al inicio del programa. La \textbf{serialización} convierte el objeto Python
completo —vectorizador y clasificador juntos— en un archivo binario
\texttt{.pkl}:

\begin{equation*}
  \text{Pipeline}(\text{vectorizador},\, \text{clasificador})
  \xrightarrow{\texttt{joblib.dump}}
  \texttt{modelo.pkl}
  \xrightarrow{\texttt{joblib.load}}
  \text{Pipeline listo para inferencia}
\end{equation*}

Esta estrategia es técnicamente relevante por tres razones:

\begin{enumerate}
  \item \textbf{Integridad:} Vectorizador y clasificador se serializan juntos,
        eliminando el riesgo de usar un clasificador entrenado con un
        vectorizador distinto al de producción.
  \item \textbf{Eficiencia:} \texttt{joblib} emplea compresión paralela y
        mapeo de memoria (\textit{memory mapping}), resultando significativamente
        más rápido que \texttt{pickle} estándar para arrays numéricos de gran
        tamaño como los parámetros internos de Sentence-BERT.
  \item \textbf{Reproducibilidad:} El archivo \texttt{.pkl} incluye la semilla
        aleatoria (\texttt{random\_state=42}) y todos los hiperparámetros,
        garantizando que las predicciones son idénticas entre sesiones y
        entornos de ejecución.
\end{enumerate}

La convención de nomenclatura de los archivos serializados es:

\begin{center}
  \texttt{<vectorizador>\_\_<clasificador>\_\_<timestamp>.pkl}
\end{center}

lo que permite trazabilidad temporal de los modelos entrenados y facilita la
comparación y auditoría de versiones en el contexto institucional de la DDP.